\documentclass[12pt]{article}

\usepackage[canadien]{babel}
\usepackage[utf8]{inputenc}
\usepackage[T1]{fontenc}

\usepackage{fancyhdr}
\usepackage{graphicx}
\usepackage{amsmath}
\usepackage{amssymb}
\usepackage{bm}
\usepackage{booktabs}
\usepackage{enumitem}
\usepackage[dvipsnames]{xcolor}
\usepackage{hyperref}
\usepackage{amsthm}

\setlength{\parindent}{0cm}
\addtolength{\oddsidemargin}{-2cm}
\addtolength{\evensidemargin}{-2cm}
\setlength{\textwidth}{17.78cm}
\addtolength{\topmargin}{-2.25cm}
\setlength{\textheight}{23.24cm}
\addtolength{\parskip}{5mm}

\pagestyle{fancy}

% --------------------
% Exercise / solution
% --------------------

\newif\ifexercise
\exercisetrue

\newif\ifsolution
\solutionfalse
% Change to \solutiontrue for the instructor version.

\theoremstyle{definition}
\newtheorem{exercise}{Question}
\newtheorem{answer}{Réponse}

\newcommand{\Exercise}[1]{%
    \ifexercise
        #1
    \fi
}

\definecolor{answercolor}{rgb}{0.00,0.12,0.60}
\newcommand{\Answer}[1]{%
    \ifsolution
        \color{answercolor}
        \begin{answer}
        #1
        \end{answer}
        \color{black}
    \fi
}

\setlist[enumerate,1]{leftmargin=*, label=\arabic*.}

\begin{document}

\fancyhead{}
\fancyfoot{}

\fancyhead[L]{
    \begin{tabular}[b]{l}
        IFT6135-A2026 \\
        Prof: Aishwarya Agrawal
    \end{tabular}
}

\fancyhead[R]{
    \begin{tabular}[b]{r}
        Devoir 1 -- Programmation \\
        Perceptrons multicouches et réseaux convolutionnels
    \end{tabular}
}

\fancyfoot[C]{- Ne pas distribuer -}

\vspace{1cm}

\textbf{Date de remise : 13 octobre 2026 à 23 h 00}

\vspace{0.2cm}

\underline{\textbf{Instructions}}

\begin{itemize}
\item \emph{Ce devoir est exigeant, commencez-le bien à l'avance !}
\item \emph{Pour toutes les questions, montrez votre démarche.}
\item \emph{La partie pratique vous a été fournie sous la forme d'un notebook Jupyter sur Google Colab ainsi que des modules  Python à compléter. Vous devez compléter les modules Python \texttt{.py} et le soumettre sur Gradescope pour l'évaluation automatique.}
\item \emph{Vous devez également soumettre votre rapport (PDF) sur Gradescope. Votre rapport doit contenir les réponses aux Questions 1.2.1 à 1.2.3 ainsi qu'aux Questions 2.1.1 à 2.2.4. Vous devez également fournir votre notebook Google Colab avec les sorties déjà générées et accessible publiquement.}
\item \emph{Les auxiliaires d'enseignement pour ce devoir sont \textbf{Achille Sowa} et \textbf{Guillermo Martin}.}
\end{itemize}
\vspace{0.2cm}

\vspace{0.4cm}


% ============================================================
% Q1
% Q1
\begin{exercise}[20] \textbf{(Implémentation d'un MLP avec NumPy)}
\Exercise{\label{ex:mlp}
\noindent Cette question est divisée en deux parties. Dans la partie~1, vous implémenterez un
perceptron multicouche (MLP) en utilisant uniquement \texttt{numpy}, ainsi que les éléments de
base nécessaires pour le définir (couches, fonctions d'activation, fonctions de perte) et
l'entraîner (optimiseur). Dans la partie~2, vous utiliserez les modèles implémentés dans la
partie~1 pour effectuer de la prévision de séries temporelles sur un ensemble de jeux de
données synthétiques.

\textbf{Lien vers le code source :}
\url{https://drive.google.com/file/d/1WxmlqjqMo1PwlZoV_-pojZGSVUZ2wrd9/view}

Quatre fichiers sources vous sont fournis à compléter~:
\begin{itemize}[nosep]
    \item \texttt{nn.py} --- définit les blocs de base. Chaque bloc de ce
    type est une sous-classe de la classe de base \texttt{Module}, et doit implémenter une
    méthode \texttt{forward} et une méthode \texttt{backward}.
    \item \texttt{model.py} --- assemble les éléments de \texttt{nn.py} pour former des modèles
    complets.
    \item \texttt{optim.py} --- implémente l'optimiseur utilisé pour entraîner les modèles (SGD).
    \item \texttt{train.py} --- entraîne les différents modèles.
\end{itemize}

\noindent \textbf{Important~:} Ne modifiez pas le code fourni, \emph{sauf} à l'intérieur des
zones délimitées par les balises \texttt{\#\#CODE} ... \texttt{\#\#END}.
Tout le reste du code doit rester inchangé~: il est requis pour la correction automatique.




\subsection*{Partie 1~: Eléments de base d'un réseau de neurones}


\begin{enumerate}[label=\textbf{1.1.\arabic*}, leftmargin=*]

    \item \textbf{Couche linéaire (2 points).} Implémentez la classe \texttt{Linear} dans
    \texttt{nn.py}, calculant
    \[
        Z = W X + b
    \]

    \begin{enumerate}[label=(\alph*), nosep]
        \item Implémentez la méthode \texttt{forward}.
        \item Implémentez la méthode \texttt{backward}, qui retourne le gradient de la perte
        par rapport à l'entrée de la couche, et qui stocke les gradients par rapport à $W$
        et $b$.
    \end{enumerate}

    \item \textbf{Fonctions d'activation (3 points).} Implémentez les trois fonctions
    d'activation suivantes en tant que sous-classes de \texttt{Module} dans \texttt{nn.py}~:
    \begin{itemize}[nosep]
        \item \texttt{ReLU}
        \item \texttt{Sigmoid}
        \item \texttt{Tanh}
    \end{itemize}
    Pour chaque classe, implémentez les méthodes  \texttt{forward} et 
    \texttt{backward}.

    \item \textbf{Fonctions de perte (2 points).} Implémentez les deux fonctions de perte
    suivantes en tant que sous-classes de \texttt{Module} dans \texttt{nn.py}~:
    \begin{itemize}[nosep]
        \item \texttt{MSE} --- erreur quadratique moyenne.
        \item \texttt{MAE} --- erreur absolue moyenne.
    \end{itemize}
    Pour chaque classe, implémentez à la fois le passage \texttt{forward} et le passage
    \texttt{backward}.

    \item \textbf{Assemblage des modèles (2 points).} À l'aide des éléments ci-dessus,
    complétez \texttt{model.py} pour construire~:
    \begin{enumerate}[label=(\alph*), nosep]
        \item un MLP à 1 couche (une seule couche linéaire, sans couche cachée)~;
        \item un MLP à 2 couches dont la fonction d'activation cachée est un argument
        configurable du constructeur (\texttt{ReLU}, \texttt{Sigmoid}, ou \texttt{Tanh}).
    \end{enumerate}

    \item \textbf{Module d'optimisation (0,5 point).} Implémentez la méthode \texttt{SGD.step}
    dans \texttt{optim.py} (voir le commentaire \texttt{TODO} dans le code pour la règle de
    mise à jour exacte).

    \item \textbf{Boucle d'entraînement (1,5 point).} Implémentez la boucle interne de la fonction \texttt{train\_model} dans \texttt{train.py} (voir le commentaire
    \texttt{TODO} dans le code pour les étapes requises).
\end{enumerate}



\subsection*{Partie 2~: Prévision de séries temporelles}

Dans cette partie, vous entraînerez les réseaux définis dans la partie~1 sur la tâche de
prévision de séries temporelles.

\subsection*{Contexte}

Une série temporelle est une séquence d'observations enregistrées à intervalles réguliers sur
une période donnée. Si une seule valeur est enregistrée à chaque intervalle, la série est
dite \emph{univariée}~; si plusieurs valeurs sont enregistrées à chaque intervalle, elle est
dite \emph{multivariée}. En recueillant $d$ valeurs à chacun des $N$ intervalles, on obtient
une série temporelle $X \in \mathbb{R}^{d \times N}$.

Étant donné un \emph{contexte} de longueur $h$, la tâche de prévision de séries temporelles
consiste à prédire les $p$ valeurs suivantes $X_{t+1:t+p}$ à partir des $h$ valeurs
précédentes $X_{t-h:t}$.

\subsection*{Jeux de données}

Un ensemble de jeux (Dataset) de données synthétiques univariés, ainsi qu'un chargeur de données (DataLoader)
correspondant, vous sont fournis afin que vous puissiez les instancier facilement. Chaque jeu
de données contient $11\,000$ points, répartis en ensembles d'entraînement, de validation et
de test selon un ratio de $6:2:2$. Les séries de base sont~:
\begin{itemize}[nosep]
    \item \texttt{sine\_slow}, \texttt{sine\_fast}
    \item \texttt{linear\_trend}
    \item \texttt{quadra\_trend}
    \item \texttt{expo\_trend}
\end{itemize}
ainsi qu'une version bruitée de chacune (bruit gaussien additif à différents rapports
signal/bruit), pour un total de $20$ jeux de données.

\subsection*{Modèles}

Instanciez les cinq modèles suivants, tous construits à partir des éléments complétées dans la
partie~1~:

\begin{center}
\begin{tabular}{@{}llll@{}}
\toprule
\textbf{Nom} & \textbf{Architecture} & \textbf{Activation cachée} & \textbf{Perte} \\
\midrule
\texttt{MLP\_1}              & MLP à 1 couche & --- & MSE \\
\texttt{MLP\_2\_ReLU}        & MLP à 2 couches & ReLU & MSE \\
\texttt{MLP\_2\_Sigmoid\_MSE} & MLP à 2 couches & Sigmoid & MSE \\
\texttt{MLP\_2\_Sigmoid\_MAE} & MLP à 2 couches & Sigmoid & MAE \\
\texttt{MLP\_2\_Tanh}        & MLP à 2 couches & Tanh & MSE \\
\bottomrule
\end{tabular}
\end{center}

\noindent Notez que \texttt{MLP\_2\_Sigmoid\_MSE} et \texttt{MLP\_2\_Sigmoid\_MAE} partagent la
même architecture et ne diffèrent que par la fonction de perte utilisée pour les entraîner, ce
qui permet une comparaison directe entre MSE et MAE.

\subsection*{Expériences}

Entraînez chacun des cinq modèles ci-dessus sur chacun des $20$ jeux de données, pour chacune
des configurations (longueur de contexte, horizon de prédiction) suivantes~:

\begin{center}
\begin{tabular}{@{}ccc@{}}
\toprule
\textbf{Configuration} & \textbf{Longueur de contexte $h$} & \textbf{Horizon de prédiction $p$} \\
\midrule
1 & 96  & 336 \\
2 & 336 & 336 \\
3 & 336 & 720 \\
\bottomrule
\end{tabular}
\end{center}

Le nombre d'époques, le taux d'apprentissage et les autres hyperparamètres d'entraînement sont
fixés à leurs valeurs par défaut directement dans \texttt{train.py}. Aucune modification de
ces paramètres n'est requise.
Il suffit d'executer le module \texttt{train.py}

\subsection*{Évaluation et analyse}

\begin{enumerate}[label=\textbf{1.2.\arabic*}, leftmargin=*]

\item \textbf{(3 points)} Pour chaque combinaison de jeu de données, de modèle et de
configuration (contexte, horizon), rapportez le \texttt{MSE} et le \texttt{MAE} de test.

Pour les jeux de données bruités, vous avez accès à la fois aux cibles bruitées et à la
vérité terrain non bruitée sous-jacente. En plus des métriques ci-dessus, calculez
\texttt{gt\_test\_mse} et \texttt{gt\_test\_mae}~: le MSE et le MAE des prédictions de chaque
modèle évalués par rapport à la vérité terrain \emph{non bruitée}, plutôt que par rapport aux
cibles bruitées. Ces métriques mesurent la capacité d'un modèle à retrouver le signal
sous-jacent réel, c'est-à-dire sa qualité en tant que débruiteur.

\item \textbf{(3 points)} À l'aide de \texttt{gt\_test\_mse} et \texttt{gt\_test\_mae},
identifiez les modèles les plus performants et les moins performants sur chaque jeu de
données, et tracez leurs prédictions. Ces résultats sont-ils cohérents avec ce à quoi vous
vous attendriez en fonction de la profondeur et du choix de la fonction d'activation de chaque
modèle~?

\item \textbf{(3 points)}
Pour les tâches ayant un horizon de prédiction $p=336$, deux longueurs de contexte sont
disponibles~: 96 et 336. Rapportez le nombre de cas où chaque contexte donne la meilleure
prédiction. Ces résultats sont-ils cohérents avec vos attentes initiales~?

\end{enumerate}

\textbf{Remise.}
Soumettez \texttt{nn.py, model.py, optim.py, train.py},
ainsi qu'un rapport PDF contenant l'analyse demandée. Le travail expérimental est évalué sur
la base de sa justesse, de sa reproductibilité et de son interprétation, et non d'une
précision cible.

\end{exercise}}




% ============================================================

% Q2
\begin{exercise}[20] \textbf{(Implémentation de couches CNN avec NumPy)}
\Exercise{\label{ex:cnn}

Implémentez des couches de convolution et de max-pooling, puis étudiez
si un CNN peut reconnaître un motif local à des positions non vues.
Cette question est indépendante de la Question~1.

Complétez le notebook fourni, qui contient les gabarits de code,
le générateur de données, le code d'entraînement et d'évaluation,
ainsi que les tests publics.
Des tests privés supplémentaires seront utilisés pour l'évaluation.

Utilisez uniquement NumPy pour vos implémentations.

\textbf{Lien vers le notebook :}
\url{https://colab.research.google.com/drive/1W7q9f-DGgVm26YqKWFL6lcwdVhOC46h7}

\subsection*{Partie 1 : Couches CNN (11 points)}

Tous les tenseurs spatiaux utilisent l'ordre $(N,C,H,W)$. Suivez les
spécifications du notebook concernant les types de données, les tailles
d'entrée valides, les caches et les gradients.

\begin{enumerate}[label=\textbf{2.1.\arabic*}, leftmargin=*]
    \item \textbf{Convolution (6 points).}
    Implémentez \texttt{conv2d\_forward} et
    \texttt{conv2d\_backward}.
    Utilisez des noyaux $3\times3$, un stride de $1$ et un pixel de
    zero-padding de chaque côté. Prenez en charge plusieurs canaux
    d'entrée et de sortie, avec un biais par canal de sortie.
    Utilisez la corrélation croisée : ne retournez pas les filtres.
    Retournez la sortie et le cache lors de la passe forward,
    ainsi que les gradients par rapport à l'entrée, aux filtres et
    aux biais lors de la passe backward.

    \item \textbf{Max pooling (3 points).}
    Implémentez \texttt{maxpool2d\_forward} et
    \texttt{maxpool2d\_backward} avec des fenêtres $2\times2$,
    un stride de $2$ et sans padding. Retournez la sortie et le cache
    lors de la passe forward, ainsi que le gradient par rapport à
    l'entrée lors de la passe backward. En cas d'égalité entre plusieurs
    maxima, sélectionnez le premier selon l'ordre ligne par ligne et
    propagez le gradient uniquement vers cette entrée.

    \item \textbf{Dimensions et paramètres (1 point).}
    Vérifiez les dimensions de sortie indiquées dans le tableau
    ci-dessous. Calculez le nombre de paramètres entraînables de chaque
    couche, biais inclus, ainsi que le nombre total de paramètres.
    Identifiez les couches qui ne possèdent aucun paramètre entraînable.

    \item \textbf{Vérifications numériques (1 point).}
    Exécutez les tests fournis ainsi que les vérifications de gradient
    par différences finies. Rapportez et sauvegardez les erreurs
    numériques. Les différences finies servent uniquement à la
    validation ; l'entraînement doit utiliser vos implémentations
    analytiques de la passe backward.
\end{enumerate}

% \newpage
\subsection*{Partie 2 : Généralisation à de nouvelles positions (9 points)}

\textbf{Jeu de données.}
Chaque image en niveaux de gris de taille $24\times24$ contient une
grille de $6\times6$ tuiles de $4\times4$ pixels : 18 tuiles sombres
d'intensité $0.2$ et 18 tuiles claires d'intensité $0.8$. Tous les
pixels d'une même tuile ont la même valeur.
Les images positives (étiquette $1$) contiennent exactement un bloc
$2\times2$ de tuiles voisines qui sont toutes sombres ou toutes
claires. Les images négatives (étiquette $0$) ne contiennent aucun
bloc de ce type. Les blocs ne dépassent pas les limites de l'image.

Les partitions fournies contiennent 2048 images d'entraînement et
512 images chacune pour la validation, le test avec positions vues
et le test avec positions non vues.
Toutes les partitions sont équilibrées entre les classes et contiennent
des arrangements de tuiles distincts.
Les exemples positifs sont également équilibrés selon la position et
le type du bloc.
Une position correspond à la coordonnée
$(\mathrm{ligne},\mathrm{colonne})$, indexée à partir de zéro, de la
tuile supérieure gauche du bloc. Les exemples positifs d'entraînement,
de validation et du test avec positions vues utilisent $(1,1)$ ou
$(1,3)$ ; les exemples positifs du test avec positions non vues
utilisent $(3,1)$ ou $(3,3)$. Les exemples négatifs suivent la même
règle d'absence de bloc dans toutes les partitions. Les jeux de test
contiennent des images différentes et ne correspondent pas à des
translations appariées.
Seules les images, et non les métadonnées décrivant l'arrangement,
sont fournies au modèle.


\begin{enumerate}[label=\textbf{2.2.\arabic*}, leftmargin=*]
    \item \textbf{Construction du CNN (2 points).}
    Complétez \texttt{build\_model(seed)} afin de retourner un nouveau
    modèle \texttt{Sequential} avec l'architecture ci-dessous.
    Utilisez l'initialisation fournie. Les deux convolutions et la
    couche linéaire incluent des biais. Les dimensions de sortie
    n'incluent pas l'axe du batch ; le modèle retourne deux logits
    bruts, sans softmax final.

    \begin{center}
    \begin{tabular}{lll}
    \toprule
    Couche & Configuration & Dimension de sortie \\
    \midrule
    Entrée & Niveaux de gris & $1\times24\times24$ \\
    Convolution + ReLU & 8 filtres, $3\times3$, stride 1, padding 1
        & $8\times24\times24$ \\
    Max pooling & $2\times2$, stride 2, sans padding
        & $8\times12\times12$ \\
    Convolution + ReLU & 16 filtres, $3\times3$, stride 1, padding 1
        & $16\times12\times12$ \\
    Global max pooling & Maximum sur les deux axes spatiaux & $16$ \\
    Linéaire & $16\rightarrow2$, avec biais & $2$ \\
    \bottomrule
    \end{tabular}
    \end{center}

    \item \textbf{Entraînement et positions vues (3 points).}
    Entraînez le modèle une seule fois pendant 30 époques avec une
    taille de batch de 64 et Adam avec un taux d'apprentissage de
    $0.003$. Conservez toutes les seeds, l'initialisation et les autres
    paramètres tels qu'ils sont fournis dans le notebook ; n'ajoutez
    aucune couche, augmentation de données ou recherche
    d'hyperparamètres.
    Sélectionnez le checkpoint ayant la plus faible cross-entropy de
    validation, en choisissant l'époque la plus ancienne en cas
    d'égalité. N'utilisez aucun des deux jeux de test pour la sélection.
    Tracez la loss et l'accuracy d'entraînement et de validation,
    rapportez l'époque sélectionnée et évaluez l'accuracy sur le jeu
    de test avec positions vues.

    \item \textbf{Positions non vues : sans réentraînement (2 points).}
    Évaluez le même checkpoint sur le jeu de test avec positions non
    vues sans mettre à jour les paramètres. Rapportez l'accuracy, le
    rappel de la classe positive et une matrice de confusion pour
    chacun des deux jeux de test, en utilisant les fonctions
    d'évaluation fournies. Utilisez les vraies étiquettes comme lignes
    de la matrice et les prédictions comme colonnes.

    \item \textbf{Interprétation (2 points).}
    Montrez au moins deux exemples positifs de chaque jeu de test,
    avec leurs étiquettes, prédictions, positions du bloc et
    identifiants d'exemple.
    Discutez d'une erreur si une erreur se produit ; sinon, indiquez
    qu'aucune erreur n'a été observée.
    Expliquez la comparaison entre positions vues et non vues en
    termes de champs réceptifs locaux, de partage des poids et de
    global max pooling. Distinguez la généralisation empirique d'une
    garantie d'invariance par translation, en tenant compte du
    zero-padding, du pooling avec stride 2 et du fait que les deux
    jeux de test contiennent des images échantillonnées indépendamment.
    Ne supposez pas que les performances doivent diminuer.
\end{enumerate}

\textbf{Remise.}
Soumettez \texttt{cnn\_layers.py}, \texttt{cnn.py}, votre notebook
exécuté, un rapport PDF contenant les calculs demandés et les résultats
expérimentaux, ainsi que l'archive exportée
\texttt{q2\_results.zip}. La partie expérimentale est évaluée selon
la justesse, la reproductibilité et l'interprétation, et non selon
une précision cible.


\Answer{
% Écrire votre réponse ici
}
\end{exercise}}

\end{document}